%%% Local Variables:
%%% mode: latex
%%% TeX-master: t
%%% End:

\documentclass[master, openany]{tongjithesis}
%开始用了openright，但这样即每章都是奇数页开始，但实际上，只需要摘要、目录、英文摘要、正文第一章、
%参考文献、附录、个人简历都从奇数页开始即可，正文内部章节可以从偶数页开始，因此用\cleardoublepage命令实现

% \documentclass[%
%   master|doctor, % mandatory option
%   xetex|pdftex|dvips|dvipdfm, % optional
%   secret,
%   openany|openright,
%   arialtoc,arialtitle]{tongjithesis}

% 所有其它可能用到的包都统一放到这里了，可以根据自己的实际添加或者删除。
\usepackage{tongjiutils}
\usepackage[top=1.3in,bottom=1.15in,left=1.25in,right=1.25in]{geometry}%调节每一页的页面大小，
\usepackage[figuresright]{rotating}
\usepackage{graphicx}
\usepackage{titlesec}
\usepackage[perpage,symbol*,bottom]{footmisc} %使用了\raggedbottom后仍然保持脚注在页面底部，且脚注在每一页都从1开始，
%符合规范，脚注形式为①，另外，此命令与下面命令还可以使得模板原先的自动填充的空白得到很好的控制，如果没有这个命令，
%章节标题与正文间的空白很难控制，特别是有图、表与公式的时候
\raggedbottom
\usepackage{epstopdf}
\usepackage{fancyhdr}
\usepackage{verbatim}
\usepackage{CJKpunct}
\usepackage{setspace}
\usepackage{longtable}
\usepackage{url}
\usepackage{pifont}
\setlength\parskip{0pt}
\frenchspacing
\widowpenalty=10000
%此命令貌似是用来修改图标与正文的间距的


% 你可以在这里修改配置文件中的定义，导言区可以使用中文。
\def\myname{同济人}

\begin{document}

% 定义所有的eps文件在 figures 子目录下
\graphicspath{{figures/}}


%%% 封面部分
\frontmatter

%%% Local Variables:
%%% mode: latex
%%% TeX-master: t
%%% End:
%\secretlevel{保密} \secretyear{2}

\ctitle{你的标题}

% 按照申请工学学位设计。如有其它需要，请修改相应文字。
\makeatletter
  \iftongji@doctor
    \cdegree{工学博士}
  \else
    \iftongji@master
      \cdegree{工学硕士}
    \fi
  \fi

\makeatother

\cdepartment{电子与信息工程学院}

\snumber{10XXXXXXXXX}

\cmajorfirst{计算机科学与技术}

\cmajorsecond{计算机应用技术}

\cauthor{你的姓名}

\csupervisor{你的 教授}

% 如果没有副指导老师或者联合指导老师，把各自{}中内容留空即可。

\cassosupervisor{}

\ccosupervisor{}

% 日期自动生成，如果你要自己写就改这个cdate
%\cdate{\CJKdigits{\the\year}年\CJKnumber{\the\month}月}
\makeatletter
  \iftongji@doctor
    \edegree{Doctor of Philosophy}
  \else
    \iftongji@master
      \edegree{Master of Science}
    \fi
  \fi

\makeatother

\cfunds{（国家自然科学基金 (No.XXXXXXXX) 支持）}

\efunds{(Supported by the Natural Science Foundation of China for\\ Grant No.XXXXXXXXX)}

\etitle{Your title}

\edepartment{School of Electrical and Informational Engineering}

\emajorfirst{Computer Science and Technology}

\emajorsecond{Computer Application Technology}

\eauthor{Tongji Ren}

\esupervisor{Prof. XXXXXX}

\eassosupervisor{}

% 这个日期也会自动生成，你要改么？
% \edate{May, 2009}

% 定义中英文摘要和关键字
\begin{cabstract}
树木建模一直是计算机图形学中一个极具挑战并且非常重要的研究课题。随着目前WebVR、WebGame、WebGIS等基于Web的应用
发展迅速，为了适应网络的传输以及用户日益增长的对图形效果的追求，如何使树木建模轻量化而富有真实感就变得尤为重要。
传统的3DSMAX、Maya等建模工具不仅耗费人力物力，而且输出的面片模型也体积庞大，不适合应用到Web领域。而诸如L-System
等基于规则的树木建模又由于其规则性而使树木模型缺失了真实感，这又不能满足用户对效果的需求。怎样在真实感和轻量化之间
进行权衡的问题亟待解决。

为了解决这个矛盾，本课题提出了一套高效、低成本，的分级轻量化树木建模方法。这里的分级轻量化体现为其对应用的适应性。
即可以基于不同应用对轻量化的不同要求，在尽可能保证真实感的前提下进行轻量化，以产生最终符合要求的模型尺寸。这种可分级
的轻量化树木建模方法还可以进一步被扩展为自动适应网络带宽条件或用户硬件条件而自动产生最匹配的轻量化树木模型的方法。

为了实现高效的分级轻量化建模方法，本文首先将PyrLK光流法进行基于仿射变换和反向追踪的改进，并且将其运用到三维重建
的特征点匹配步骤中，以提高树木特征点的匹配率和稳定性。然后进行GPU加速的三维重建以得到高精度点云模型。接着本文
运用三维体素泛洪和最小二乘线性拟合的方法对树木骨架和半径信息进行抽取，以适应树木生长规律的方法抽取出了准确的骨架。
然后本文提出了根据应用对轻量化的需求等级，对骨架进行纵向和横向的合并，以减小骨架的尺寸来实现轻量化，从而更好地适应
面向网络的应用的需求。最后本文还给出了一套完整的基于图像树木建模的质量评价，提出了还原度的概念来客观、量化地评价建模出
的模型的还原度以及在轻量化过程中质量的丢失。
\end{cabstract}

\ckeywords{基于图像建模， 树木建模， 轻量化建模， 三维重建， 骨架抽取}

\begin{eabstract}
Tree modeling has long been a challenging subject in computer graphics. As the Web-oriented applications(WebVR, WebGame,
WebGIS, etc) develop rapidly and the persuit of graphics effect increases quickly, the lightweight and realism of tree modeling are badly
needed nowadays. The traditional 3d modeling tools such as 3DSMAX and Maya are not only time and labour consuming, but it
takes a large model size, which are not practical to Web apps. The rule-based modeling such as L-System can solve the size
problem, but its output lacks realism, which is not tolerated by users. So the balance between realism and lightweight
is a real problem which are eagerly demanded to solve.

In order to solve this problem, a high-efficiency, low-cost and lightweight-classified tree modeling method is proposed.
Here the lightweight-classified means it can produce different lightweight levels of tree models. And to implement this 
goal, this method will reduce model size on the premise of not losing much realism. This method can also be applied furthur
to automatically adapt to the bandwidth and hardware conditions of the client side.

For implementing the lightweight method, we first improve the traditional PyrLK optical flow method to support affine transformation
and backward feature tracking, which can furthur be applied to do feature matching in gpu accelerated 3D reconstruction and 
improve the match ratio. Then we use flooding algorithm in 3D voxel model and least squares method to discover the tree skeleton and
its radius information. According to the lightweight level the applications require, we reduce the model size by merging the
branches vertically and horizontally respectively. At last we propose a modeling quality evaluation method, which will objectively and 
quantizedly evaluate the restore degree of the tree model.
\end{eabstract}

\ekeywords{Image-based Modeling, Tree Modeling, Light-weight Modeling, PyrLK Optical Flow, Skeleton Extraction}

\makecover


% 目录

\tableofcontents

% 符号对照表
\begin{denotation}
\item[GNU] GNU's Not Unix /'gnu:/
\item[GFDL] GNU Free Documentation License
\item[GPL] GNU General Public License
\item[FSF] Free Software Foundation
\item[SMP] 对称多处理
\item[API] 应用程序编程接口
\item[$E$] 能量
\item[$m$] 质量
\item[$c$] 光速
\item[$P$] 概率
\item[$T$] 时间
\item[$v$] 速度
\end{denotation}


%%% 以下索引按需要选择
% 插图索引
%\listoffigures
% 表格索引
%\listoftables
% 公式索引
%\listofequations

%%% 正文
\mainmatter

%%% Local Variables:
%%% mode: latex
%%% TeX-master: t
%%% End:



\chapter{引言}
\label{cha:intro}
\section{背景介绍}
\label{sec:background}
在互联网飞速发展的今天，网络应用已经延伸到生活的方方面面。微博、人人网、
在线购物、在线音乐等已经成为当今人们生活的一部分。面向Web的虚拟现实应用
如WebVR、WebGame、WebGIS等也必然将成为虚拟现实发展的重要方向。树木作为自然
界常见的事物，在各种虚拟现实的场景中出现的频率很高。然而树木形态各异，结构
复杂，给3D建模带来了很大的难度。通常单棵树的数据量已经不小，对于构建一个
树木的聚集场景(如森林)就更加庞大，这容易使得场景负荷变大而产生延迟。因此，
树木建模的质量和效率将直接决定面向Web的虚拟现实应用的成败。

目前的树木的3D建模，主要是通过专业的3D建模工具(3DSMAX、Maya等)进行手工建模。
这种建模方法对建模人员的要求较高，并且需要的时间长。而且这种方法通常最终生成
的是面片信息，要表达一棵形态复杂的树木需要大量的顶点信息，导致最终生成的模型
体积较大，对于需要大批树木的场景，负荷就会变得更大。

目前树木的轻量化建模，从最简单的基于分形，广告牌技术的建模到稍微复杂的基于
规则的建模，都存在一个共同的问题，就是在轻量化的同时，很大程度上舍弃了模型
的真实感和树木本身的形态特征。随着当今应用对真实度要求的升高，这类轻量化的
建模方法已经不能完全满足需求。真实感与轻量化之间的权衡也成为了当今应用需要
考虑的一个重要因素。

本课题基于以上的考虑，从基于图片对树木结构进行完整的恢复，到面向应用需要对
真实感与轻量化进行人工控制，到最后模型重建质量的评估，给出了一套完整的解决
方案。

\section{课题的主要工作}
\label{sec:objective}
本课题的主要工作有：

1. 对PyrLK光流法进行改良，并将其运用于三维重建算法中的特征点匹配步骤，使树木
重建的模型更加准确和精细。

2. 提出了基于三维体素泛洪和空间反向投影的三维重建方法，以连续的体素替代传统
方法不连续的点云，使得后续的骨架抽取步骤更加准确和方便。

3. 提出了基于多方向迭代和步长探索的树木骨架抽取方法。该方法区别于传统的3D瘦化
骨架抽取方法，它只适用于具有分形结构的3D骨架，所以更能够得到准确的树木骨架。

4. 提出了基于用户交互对树木模型进行完善和轻量化，让最终的应用来决定其所需的树木
模型，避免了主观的一味轻量化或一味追求真实感而带来的需求矛盾，将模型的成型延迟
至具体的应用。

5. 提出了基于图像的树木重建质量评估方法，对于建模质量和轻量化过程中真实感的下降
程度给出了函数化和量化的评价依据。


%%% Local Variables:
%%% mode: latex
%%% TeX-master: t
%%% End:



\chapter{中华人民共和国}
\label{cha:china}

\section{其它例子}
\label{sec:other}

在第~\ref{cha:intro} 章中我们学习了贝叶斯公式~(\ref{equ:chap1:bayes})，这里我们复
习一下：
\begin{equation}
\label{equ:chap2:bayes}
p(y|\mathbf{x}) = \frac{p(\mathbf{x},y)}{p(\mathbf{x})}=
\frac{p(\mathbf{x}|y)p(y)}{p(\mathbf{x})}
\end{equation}

\subsection{绘图}
\label{sec:draw}

本模板不再预先装载任何绘图包（如 \textsf{pstricks，pgf} 等），完全由你自己来决定。
个人觉得 \textsf{pgf} 不错，不依赖于 Postscript。此外还有很多针对 \LaTeX{} 的
 GUI 作图工具，如 XFig(jFig), WinFig, Tpx, Ipe, Dia, Inkscape, LaTeXPiX,
jPicEdt, jaxdraw 等等。

\subsection{插图}
\label{sec:graphs}

强烈推荐《\LaTeXe 插图指南》！关于子图形的使用细节请参看 \textsf{subfig} 的说明文档。

\subsubsection{一个图形}
\label{sec:onefig}
一般图形都是处在浮动环境中。之所以称为浮动是指最终排版效果图形的位置不一定与源文
件中的位置对应\footnote{This is not a bug, but a feature of
\LaTeX!}，这也是刚使 用 \LaTeX{}
同学可能遇到的问题。如果要强制固定浮动图形的位置，请使用
\textsf{float} 宏包， 它提供了 \texttt{[H]}
参数，比如图~\ref{fig:heythere}。
\begin{figure}[H] % use float package if you want it here
  \centering
  \includegraphics[height=2cm]{hello.jpg}
  \caption{插个图插个图}
  \label{fig:heythere}
\end{figure}

大学之道，在明明德，在亲民，在止于至善。知止而后有定；定而后能静；静而后能安；安
而后能虑；虑而后能得。物有本末，事有终始。知所先后，则近道矣。古之欲明明德于天
下者，先治其国；欲治其国者，先齐其家；欲齐其家者，先修其身；欲修其身者，先正其心；
欲正其心者，先诚其意；欲诚其意者，先致其知；致知在格物。物格而后知至；知至而后
意诚；意诚而后心正；心正而后身 修；身修而后家齐；家齐而后国治；国治而后天下
平。自天子以至于庶人，壹是皆以修身为本。其本乱而未治者 否矣。其所厚者薄，而其所
薄者厚，未之有也！

\hfill \pozhehao《大学》


\subsubsection{多个图形}
\label{sec:multifig}

如果多个图形相互独立，并不共用一个图形计数器，那么用 \verb|minipage| 或者
\verb|parbox| 就可以。否则，请参看图~\ref{fig:big1}，它包含两个小图，分别是图~\ref{fig:subfig1}
和图~\ref{fig:subfig2}。推荐使用 \verb|\subfloat|，不要再用
\verb|\subfigure| 和 \verb|\subtable|。
\begin{figure} %[h]
  \centering%
  \subfloat[第一个小图形]{%
    \label{fig:subfig1}
    \includegraphics[height=2cm]{tongji-fig-logo.png}}\hspace{4em}%
  \subfloat[第二个小图形。如果标题很长的话，它会自动换行，这个 caption 就是这样的例子]{%
    \label{fig:subfig2}
    \includegraphics[height=2cm]{tongji-text-logo.png}}
  \caption{包含子图形的大图形}
  \label{fig:big1}
\end{figure}

古之学者必有师。师者，所以传道受业解惑也。人非生而知之者，孰能无惑？惑而不从师，
其为惑也，终不解矣。生乎吾前，其闻道也固先乎吾，吾从而师之；生乎吾後，其闻道也亦
先乎吾，吾从而师之。吾师道也，夫庸知其年之先後生於吾乎！是故无贵无贱无长无少，道
之所存，师之所存也。

嗟乎！师道之不传也久矣，欲人之无惑也难矣。古之圣人，其出人也远矣，犹且从师而问焉；
今之众人，其下圣人也亦远矣，而耻学於师。是故圣益圣，愚益愚。圣人之所以为圣，愚
人之所以为愚，其皆出於此乎？爱其子，择师而教之，於其身也，则耻师焉，惑焉。彼童子
之师，授之书而习其句读者，非吾所谓传其道、解其惑者也。句读之不知，惑之不解，或师
焉，或不焉，小学而大遗，吾未见其明也。巫医、乐师、百工之人不耻相师，  士大夫之族
曰“师”曰“弟子”之云者，则群聚而笑之。问之，则曰：彼与彼年相若也，道相似也，位
卑则足羞，官盛则近谀。呜呼！师道之不复，可知矣。巫医、乐师、百工之人。吾子不齿，
今其智乃反不能及，其可怪也欤！圣人无常师。孔子师郯子、苌子、师襄、老聃。郯子之徒，
其贤不及孔子。孔子曰：“三人行，必有我师。”是故弟子不必不如师，师不必贤於弟子。
闻道有先後，术业有专攻，如是而已。

如果要把编号的两个图形并排，那么小页就非常有用了：
\begin{figure}
\begin{minipage}{0.48\textwidth}
  \centering
  \includegraphics[height=2cm]{tongji-whole-logo.png}
  \caption{并排第一个图}
  \label{fig:parallel1}
\end{minipage}\hfill
\begin{minipage}{0.48\textwidth}
  \centering
  \includegraphics[height=2cm]{tongji-whole-logo.png}
  \caption{并排第二个图}
  \label{fig:parallel2}
\end{minipage}
\end{figure}

李氏子蟠，年十七，好古文、六艺，经传皆通习之，不拘於时，学於余。余嘉其能行古
道，作师说以贻之。

\hfill \pozhehao 韩愈（唐）


%%% 其它部分
\backmatter

\makeatother

% 致谢
\cleardoublepage

%%% Local Variables:
%%% mode: latex
%%% TeX-master: "../main"
%%% End:

\begin{ack}
  衷心感谢导师 xxx 教授和物理系 xxx 副教授对本人的精心指导。他们的言传身教将使
  我终生受益。

  在美国麻省理工学院化学系进行九个月的合作研究期间，承蒙 xxx 教授热心指导与帮助，不
  胜感激。感谢 xx 实验室主任 xx 教授，以及实验室全体老师和同学们的热情帮助和支
  持！本课题承蒙国家自然科学基金资助，特此致谢。

  感谢 \tongjithesis{}，它的存在让我的论文写作轻松自在了许多，让我的论文格式规整漂亮了
  许多。

\end{ack}


% 参考文献
\cleardoublepage
\bibliographystyle{tongjibib}
\bibliography{ref/refs}


% 附录
\cleardoublepage
\begin{appendix}
%%% Local Variables:
%%% mode: latex
%%% TeX-master: "../main"
%%% End:

\chapter{外文资料原文}
\label{cha:engorg}
As one of the most widely used techniques in operations research, {\em
  mathematical programming} is defined as a means of maximizing a quantity known
as {\em objective function}, subject to a set of constraints represented by
equations and inequalities. Some known subtopics of mathematical programming are
linear programming, nonlinear programming, multiobjective programming, goal
programming, dynamic programming, and multilevel programming$^{[1]}$.

It is impossible to cover in a single chapter every concept of mathematical
programming. This chapter introduces only the basic concepts and techniques of
mathematical programming such that readers gain an understanding of them
throughout the book$^{[2,3]}$.


\section{Single-Objective Programming}
The general form of single-objective programming (SOP) is written
as follows,
\begin{equation}\tag*{(123)} % 如果附录中的公式不想让它出现在公式索引中，那就请
                             % 用 \tag*{xxxx}
\left\{\begin{array}{l}
\max \,\,f(x)\\[0.1 cm]
\mbox{subject to:} \\ [0.1 cm]
\qquad g_j(x)\le 0,\quad j=1,2,\cdots,p
\end{array}\right.
\end{equation}
which maximizes a real-valued function $f$ of
$x=(x_1,x_2,\cdots,x_n)$ subject to a set of constraints.

One of the outstanding contributions to mathematical programming was known as
the Kuhn-Tucker conditions\ref{eq:ktc}. In order to introduce them, let us give
some definitions. An inequality constraint $g_j(x)\le 0$ is said to be active at
a point $x^*$ if $g_j(x^*)=0$. A point $x^*$ satisfying $g_j(x^*)\le 0$ is said
to be regular if the gradient vectors $\nabla g_j(x)$ of all active constraints
are linearly independent.

%\newtheorem{mpdef}{Definition}[chapter]
%\begin{mpdef}
%In SOP, we call $x$ a decision vector, and
%$x_1,x_2,\cdots,x_n$ decision variables. The function
%$f$ is called the objective function. The set
%\begin{equation}\tag*{(456)} % 这里同理，其它不再一一指定。
%S=\left\{x\in\Re^n\bigm|g_j(x)\le 0,\,j=1,2,\cdots,p\right\}
%\end{equation}
%is called the feasible set. An element $x$ in $S$ is called a
%feasible solution.
%\end{mpdef}
%
%\newtheorem{mpdefop}[mpdef]{Definition}
%\begin{mpdefop}
%A feasible solution $x^*$ is called the optimal
%solution of SOP if and only if
%\begin{equation}
%f(x^*)\ge f(x)
%\end{equation}
%for any feasible solution $x$.
%\end{mpdefop}

Let $x^*$ be a regular point of the constraints of SOP and assume that all the
functions $f(x)$ and $g_j(x),j=1,2,\cdots,p$ are differentiable. If $x^*$ is a
local optimal solution, then there exist Lagrange multipliers
$\lambda_j,j=1,2,\cdots,p$ such that the following Kuhn-Tucker conditions hold,
\begin{equation}
\label{eq:ktc}
\left\{\begin{array}{l}
    \nabla f(x^*)-\sum\limits_{j=1}^p\lambda_j\nabla g_j(x^*)=0\\[0.3cm]
    \lambda_jg_j(x^*)=0,\quad j=1,2,\cdots,p\\[0.2cm]
    \lambda_j\ge 0,\quad j=1,2,\cdots,p.
\end{array}\right.
\end{equation}
If all the functions $f(x)$ and $g_j(x),j=1,2,\cdots,p$ are convex and
differentiable, and the point $x^*$ satisfies the Kuhn-Tucker conditions
(\ref{eq:ktc}), then it has been proved that the point $x^*$ is a global optimal
solution of SOP.

\subsection{Linear Programming}
\label{sec:lp}

If the functions $f(x),g_j(x),j=1,2,\cdots,p$ are all linear, then SOP is called
a {\em linear programming}.

The feasible set of linear is always convex. A point $x$ is called an extreme
point of convex set $S$ if $x\in S$ and $x$ cannot be expressed as a convex
combination of two points in $S$. It has been shown that the optimal solution to
linear programming corresponds to an extreme point of its feasible set provided
that the feasible set $S$ is bounded. This fact is the basis of the {\em simplex
  algorithm} which was developed by Dantzig as a very efficient method for
solving linear programming.
\begin{table}[ht]
\centering
  \centering
  \caption*{Table~1\hskip1em This is an example for manually numbered table, which
    would not appear in the list of tables}
  \label{tab:badtabular2}
  \begin{tabular}[c]{|c|m{0.8in}|c|c|c|c|c|}\hline
    \multicolumn{2}{|c|}{Network Topology} & \# of nodes &
    \multicolumn{3}{c|}{\# of clients} & Server \\\hline
    GT-ITM & Waxman Transit-Stub & 600 &
    \multirow{2}{2em}{2\%}&
    \multirow{2}{2em}{10\%}&
    \multirow{2}{2em}{50\%}&
    \multirow{2}{1.2in}{Max. Connectivity}\\\cline{1-3}
    \multicolumn{2}{|c|}{Inet-2.1} & 6000 & & & &\\\hline
    \multirow{2}{1in}{Xue} & Rui  & Ni &\multicolumn{4}{c|}{\multirow{2}*{\tongjithesis}}\\\cline{2-3}
    & \multicolumn{2}{c|}{ABCDEF} &\multicolumn{4}{c|}{} \\\hline
\end{tabular}
\end{table}

Roughly speaking, the simplex algorithm examines only the extreme points of the
feasible set, rather than all feasible points. At first, the simplex algorithm
selects an extreme point as the initial point. The successive extreme point is
selected so as to improve the objective function value. The procedure is
repeated until no improvement in objective function value can be made. The last
extreme point is the optimal solution.

\subsection{Nonlinear Programming}

If at least one of the functions $f(x),g_j(x),j=1,2,\cdots,p$ is nonlinear, then
SOP is called a {\em nonlinear programming}.

A large number of classical optimization methods have been developed to treat
special-structural nonlinear programming based on the mathematical theory
concerned with analyzing the structure of problems.
\begin{figure}[h]
  \centering
  \includegraphics[clip]{tongji-lib-logo.jpg}
  \caption*{Figure~1\hskip1em This is an example for manually numbered figure,
    which would not appear in the list of figures}
  \label{tab:badfigure2}
\end{figure}

Now we consider a nonlinear programming which is confronted solely with
maximizing a real-valued function with domain $\Re^n$.  Whether derivatives are
available or not, the usual strategy is first to select a point in $\Re^n$ which
is thought to be the most likely place where the maximum exists. If there is no
information available on which to base such a selection, a point is chosen at
random. From this first point an attempt is made to construct a sequence of
points, each of which yields an improved objective function value over its
predecessor. The next point to be added to the sequence is chosen by analyzing
the behavior of the function at the previous points. This construction continues
until some termination criterion is met. Methods based upon this strategy are
called {\em ascent methods}, which can be classified as {\em direct methods},
{\em gradient methods}, and {\em Hessian methods} according to the information
about the behavior of objective function $f$. Direct methods require only that
the function can be evaluated at each point. Gradient methods require the
evaluation of first derivatives of $f$. Hessian methods require the evaluation
of second derivatives. In fact, there is no superior method for all
problems. The efficiency of a method is very much dependent upon the objective
function.

\subsection{Integer Programming}

{\em Integer programming} is a special mathematical programming in which all of
the variables are assumed to be only integer values. When there are not only
integer variables but also conventional continuous variables, we call it {\em
  mixed integer programming}. If all the variables are assumed either 0 or 1,
then the problem is termed a {\em zero-one programming}. Although integer
programming can be solved by an {\em exhaustive enumeration} theoretically, it
is impractical to solve realistically sized integer programming problems. The
most successful algorithm so far found to solve integer programming is called
the {\em branch-and-bound enumeration} developed by Balas (1965) and Dakin
(1965). The other technique to integer programming is the {\em cutting plane
  method} developed by Gomory (1959).

\hfill\textit{Uncertain Programming\/}\quad(\textsl{BaoDing Liu, 2006.2})

\section*{References}
\noindent{\itshape NOTE: these references are only for demonstration, they are
  not real citations in the original text.}

\begin{enumerate}[{$[$}1{$]$}]
\item Donald E. Knuth. The \TeX book. Addison-Wesley, 1984. ISBN: 0-201-13448-9
\item Paul W. Abrahams, Karl Berry and Kathryn A. Hargreaves. \TeX\ for the
  Impatient. Addison-Wesley, 1990. ISBN: 0-201-51375-7
\item David Salomon. The advanced \TeX book.  New York : Springer, 1995. ISBN:0-387-94556-3
\end{enumerate}

\chapter{外文资料的调研阅读报告或书面翻译}
\section{单目标规划}
北冥有鱼，其名为鲲。鲲之大，不知其几千里也。化而为鸟，其名为鹏。鹏之背，不知其几
千里也。怒而飞，其翼若垂天之云。是鸟也，海运则将徙于南冥。南冥者，天池也。
\begin{equation}\tag*{(123)}
 p(y|\mathbf{x}) = \frac{p(\mathbf{x},y)}{p(\mathbf{x})}=
\frac{p(\mathbf{x}|y)p(y)}{p(\mathbf{x})}
\end{equation}

吾生也有涯，而知也无涯。以有涯随无涯，殆已！已而为知者，殆而已矣！为善无近名，为
恶无近刑，缘督以为经，可以保身，可以全生，可以养亲，可以尽年。

\subsection{线性规划}
庖丁为文惠君解牛，手之所触，肩之所倚，足之所履，膝之所倚，砉然响然，奏刀騞然，莫
不中音，合于桑林之舞，乃中经首之会。
\begin{table}[ht]
\centering
  \centering
  \caption*{表~1\hskip1em 这是手动编号但不出现在索引中的一个表格例子}
  \label{tab:badtabular3}
  \begin{tabular}[c]{|c|m{0.8in}|c|c|c|c|c|}\hline
    \multicolumn{2}{|c|}{Network Topology} & \# of nodes &
    \multicolumn{3}{c|}{\# of clients} & Server \\\hline
    GT-ITM & Waxman Transit-Stub & 600 &
    \multirow{2}{2em}{2\%}&
    \multirow{2}{2em}{10\%}&
    \multirow{2}{2em}{50\%}&
    \multirow{2}{1.2in}{Max. Connectivity}\\\cline{1-3}
    \multicolumn{2}{|c|}{Inet-2.1} & 6000 & & & &\\\hline
    \multirow{2}{1in}{Xue} & Rui  & Ni &\multicolumn{4}{c|}{\multirow{2}*{\tongjithesis}}\\\cline{2-3}
    & \multicolumn{2}{c|}{ABCDEF} &\multicolumn{4}{c|}{} \\\hline
\end{tabular}
\end{table}

文惠君曰：“嘻，善哉！技盖至此乎？”庖丁释刀对曰：“臣之所好者道也，进乎技矣。始臣之
解牛之时，所见无非全牛者；三年之后，未尝见全牛也；方今之时，臣以神遇而不以目视，
官知止而神欲行。依乎天理，批大郤，导大窾，因其固然。技经肯綮之未尝，而况大坬乎！
良庖岁更刀，割也；族庖月更刀，折也；今臣之刀十九年矣，所解数千牛矣，而刀刃若新发
于硎。彼节者有间而刀刃者无厚，以无厚入有间，恢恢乎其于游刃必有余地矣。是以十九年
而刀刃若新发于硎。虽然，每至于族，吾见其难为，怵然为戒，视为止，行为迟，动刀甚微，
謋然已解，如土委地。提刀而立，为之而四顾，为之踌躇满志，善刀而藏之。”

文惠君曰：“善哉！吾闻庖丁之言，得养生焉。”


\subsection{非线性规划}
孔子与柳下季为友，柳下季之弟名曰盗跖。盗跖从卒九千人，横行天下，侵暴诸侯。穴室枢
户，驱人牛马，取人妇女。贪得忘亲，不顾父母兄弟，不祭先祖。所过之邑，大国守城，小
国入保，万民苦之。孔子谓柳下季曰：“夫为人父者，必能诏其子；为人兄者，必能教其弟。
若父不能诏其子，兄不能教其弟，则无贵父子兄弟之亲矣。今先生，世之才士也，弟为盗
跖，为天下害，而弗能教也，丘窃为先生羞之。丘请为先生往说之。”
\begin{figure}[h]
  \centering
  \includegraphics{hello}
  \caption*{图~1\hskip1em 这是手动编号但不出现索引中的图片的例子}
  \label{tab:badfigure3}
\end{figure}

柳下季曰：“先生言为人父者必能诏其子，为人兄者必能教其弟，若子不听父之诏，弟不受
兄之教，虽今先生之辩，将奈之何哉？且跖之为人也，心如涌泉，意如飘风，强足以距敌，
辩足以饰非。顺其心则喜，逆其心则怒，易辱人以言。先生必无往。”

孔子不听，颜回为驭，子贡为右，往见盗跖。

\subsection{整数规划}
盗跖乃方休卒徒大山之阳，脍人肝而餔之。孔子下车而前，见谒者曰：“鲁人孔丘，闻将军
高义，敬再拜谒者。”谒者入通。盗跖闻之大怒，目如明星，发上指冠，曰：“此夫鲁国之
巧伪人孔丘非邪？为我告之：尔作言造语，妄称文、武，冠枝木之冠，带死牛之胁，多辞缪
说，不耕而食，不织而衣，摇唇鼓舌，擅生是非，以迷天下之主，使天下学士不反其本，妄
作孝弟，而侥幸于封侯富贵者也。子之罪大极重，疾走归！不然，我将以子肝益昼餔之膳。”


\chapter{其它附录}
其它附录的内容可以放到这里，当然如果你愿意，可以把这部分也放到独立的文件中，然后
将其\verb|\input| 到主文件中。
\section{测试附录}

\end{appendix}

% 个人简历
\cleardoublepage
\begin{resume}

  \resumeitem{个人简历}

  xxxx 年 xx 月 xx 日出生于 xx 省 xx 县。
  
  xxxx 年 9 月考入 xx 大学 xx 系 xx 专业，xxxx 年 7 月本科毕业并获得 xx 学士学位。
  
  xxxx 年 9 月免试进入 xx 大学 xx 系攻读 xx 学位至今。

  \resumeitem{发表的学术论文} % 发表的和录用的合在一起

  \begin{enumerate}[{[}1{]}]
  \item Yang Y, Ren T L, Zhang L T, et al. Miniature microphone with silicon-
    based ferroelectric thin films. Integrated Ferroelectrics, 2003,
    52:229-235. (SCI 收录, 检索号:758FZ.)
  \item 杨轶, 张宁欣, 任天令, 等. 硅基铁电微声学器件中薄膜残余应力的研究. 中国机
    械工程, 2005, 16(14):1289-1291. (EI 收录, 检索号:0534931 2907.)
  \item 杨轶, 张宁欣, 任天令, 等. 集成铁电器件中的关键工艺研究. 仪器仪表学报,
    2003, 24(S4):192-193. (EI 源刊.)
  \item Yang Y, Ren T L, Zhu Y P, et al. PMUTs for handwriting recognition. In
    press. (已被 Integrated Ferroelectrics 录用. SCI 源刊.)
  \item Wu X M, Yang Y, Cai J, et al. Measurements of ferroelectric MEMS
    microphones. Integrated Ferroelectrics, 2005, 69:417-429. (SCI 收录, 检索号
    :896KM.)
  \item 贾泽, 杨轶, 陈兢, 等. 用于压电和电容微麦克风的体硅腐蚀相关研究. 压电与声
    光, 2006, 28(1):117-119. (EI 收录, 检索号:06129773469.)
  \item 伍晓明, 杨轶, 张宁欣, 等. 基于MEMS技术的集成铁电硅微麦克风. 中国集成电路, 
    2003, 53:59-61.
  \end{enumerate}

  \resumeitem{研究成果} % 有就写，没有就删除
  \begin{enumerate}[{[}1{]}]
  \item 任天令, 杨轶, 朱一平, 等. 硅基铁电微声学传感器畴极化区域控制和电极连接的
    方法: 中国, CN1602118A. (中国专利公开号.)
  \item Ren T L, Yang Y, Zhu Y P, et al. Piezoelectric micro acoustic sensor
    based on ferroelectric materials: USA, No.11/215, 102. (美国发明专利申请号.)
  \end{enumerate}
\end{resume}


\end{document}
